\section{Virtual Representation Theorem}\label{sec:injective_symmetry_cocycle}

With a full group of on-site symmetries (rather than just a single
unitary), the resulting gauges determine a projective representation
on the bond space.

\begin{definition}[Twisted tensor]\label{def:twisted_tensor}
    \lean{MPSTensor.twistedTensor}
    \leanok
    \uses{def:mps_tensor}
    Given a group representation $U : G \to \GL_{d}(\C)$ on the physical
    index, the \emph{$g$-twisted tensor} is defined, for each
    $i \in \{0,\ldots,d{-}1\}$, by
    \begin{align}
        \widetilde{A}_g^i
        &:=\sum_{j=0}^{d-1}U(g)_{ij}\,A^j.
        \label{eq:symmetry_twisted_tensor}
    \end{align}
\end{definition}

\begin{definition}[On-site symmetry]\label{def:on_site_symmetric}
    \lean{MPSTensor.IsOnSiteSymmetric}
    \leanok
    \uses{def:twisted_tensor, def:same_mpv}
    A tensor $A$ is \emph{on-site symmetric} under a group
    representation $U : G \to \GL_{d}(\C)$ if
    $\mc{V}(A) = \mc{V}(\widetilde{A}_g)$ for every $g \in G$.
\end{definition}

\begin{lemma}[Functoriality of twisting]\label{lem:twisted_tensor_mul}
    \lean{MPSTensor.twistedTensor_mul}
    \leanok
    \uses{def:twisted_tensor}
    The twisted tensor satisfies the composition law
    \begin{align}
        \widetilde{A}_{gh}
        &=\widetilde{(\widetilde{A}_h)}_g,
        \label{eq:symmetry_twist_composition}
    \end{align}
    i.e.\ twisting by $gh$ is the same as first twisting by~$h$
    and then by~$g$.
\end{lemma}

\begin{proof}\leanok\uses{def:twisted_tensor}
    Expand both sides of (\ref{eq:symmetry_twist_composition}) using
    (\ref{eq:symmetry_twisted_tensor}) and $U(gh)=U(g)\,U(h)$, then swap
    the order of summation.
\end{proof}

\begin{lemma}[Identity twist]\label{lem:twisted_tensor_one}
    \lean{MPSTensor.twistedTensor_one}
    \leanok
    \uses{def:twisted_tensor}
    Twisting by the identity group element is trivial:
    $\widetilde{A}_1=A$.
\end{lemma}

\begin{proof}\leanok\uses{def:twisted_tensor}
    Since $U(1)=\Id_d$, each component in
    (\ref{eq:symmetry_twisted_tensor}) reduces to
    $\sum_j\delta_{ij}\,A^j=A^i$.
\end{proof}

\begin{lemma}[Symmetric twists are gauge equivalent]%
    \label{lem:gauge_equiv_twisted}
    \lean{MPSTensor.gaugeEquiv_twistedTensor_of_injective}
    \leanok
    \uses{def:injective, def:on_site_symmetric}
    If $A$ is injective and on-site symmetric under $U$, then
    for every $g \in G$ the twisted tensor $\widetilde{A}_g$ is
    gauge equivalent to~$A$.
\end{lemma}

\begin{proof}\leanok\uses{def:on_site_symmetric, thm:ft_single}
    On-site symmetry gives
    $\mc{V}(A)=\mc{V}(\widetilde{A}_g)$, and the single-block
    Fundamental Theorem turns this equality of MPV families into gauge
    equivalence.
\end{proof}

\begin{remark}[Fundamental-Theorem route for injective symmetry]%
    \label{rem:injective_symmetry_ft_route}
    \uses{thm:ft_single, lem:gauge_equiv_twisted, lem:gauge_unique_up_to_scalar,
        lem:twisted_tensor_mul, thm:virtual_rep_injective}
    The injective virtual-representation theorem is not proved from string-order
    rigidity. The on-site symmetry condition first identifies $A$ and
    $\widetilde{A}_g$ as the same MPV family; the single-block Fundamental
    Theorem gives the gauge in Lemma~\ref{lem:gauge_equiv_twisted}. Gauge
    uniqueness and the twist composition law then give
    Theorem~\ref{thm:virtual_rep_injective}. The string-order results later in
    this chapter use this virtual representation as an input.
\end{remark}

\begin{theorem}[Virtual symmetry equation]\label{thm:virtual_symmetry_eq}
    \lean{MPSTensor.virtual_symmetry_eq}
    \leanok
    \uses{def:injective, def:on_site_symmetric, lem:gauge_equiv_twisted}
    If $A$ is injective and on-site symmetric under $U$, then
    for every $g \in G$ there exist an invertible matrix
    $X(g) \in \GL_{D}(\C)$ and a nonzero scalar
    $\varphi(g) \in \C^{\times}$ (in fact $\varphi=1$
    in the single-block case) such that, for every physical index~$i$,
    \begin{align}
        \widetilde{A}_g^{\,i}
        &=\varphi(g)\,X(g)\,A^i\,X(g)^{-1}.
        \label{eq:symmetry_virtual_equation}
    \end{align}
\end{theorem}

\begin{proof}\leanok\uses{lem:gauge_equiv_twisted}
    The physical action is transferred to the virtual level by the same
    local move that replaces the twisted tensor by a gauge-conjugated one:
    % Formula: (\widetilde A_g^i)_{ab}=\sum_j U(g)_{ij}(A^j)_{ab}.
    % Source verdict: U(g) acts only on the physical index of A and has
    % no virtual indices.
    % Index map: i is the free upper physical index, j is the unique
    % contracted index, and a,b are the free west/east virtual indices.
    % Boundary: virtual-west=1 | virtual-east=1 | physical-up=1 |
    % physical-down=0.
    \begin{tenkzequation}
        \begin{tenkz}[rows={op:none,ket}]
            \tn[role=operator,up=$i$]{U(g)} \\
            \tn{A}
        \end{tenkz}
        \(\qquad \leadsto \qquad\)
        % Formula/source: Papers/0802.0447/StringOrder-v10.tex,
        % lines 328--349, especially Condition C1, gives the unitary u,V
        % specialization (V^\dagger = V^{-1}); the surrounding fundamental
        % theorem supplies the analogous general GL gauge relation displayed here.
        % Ink-to-index map: i is the free upper physical index; the two
        % horizontal wires are the free virtual indices.
        % Contraction set: only X(g)--A and A--X(g)^{-1} are summed.
        % Type verdict: a rank-three twisted local tensor.
        % Boundary signature: virtual-west=1 | virtual-east=1 |
        % physical-up=1 | physical-down=0.
        \begin{tenkz}[physical=up]
            \tnX{X(g)} & \tn[up=$i$]{A} & \tnX{X(g)^{-1}}
        \end{tenkz}.
    \end{tenkzequation}
    Apply Lemma~\ref{lem:gauge_equiv_twisted} to obtain $X(g)$, then
    set $\varphi(g)=1$ in (\ref{eq:symmetry_virtual_equation}).
\end{proof}

\begin{lemma}[Gauge uniqueness up to scalar]\label{lem:gauge_unique_up_to_scalar}
    \lean{MPSTensor.gauge_unique_up_to_scalar}
    \leanok
    \uses{def:injective, def:gauge_equiv, lem:center_scalar}
    Let $A$ be an injective tensor.
    If $X,Y \in \GL_{D}(\C)$ both satisfy
    $B^i=X\,A^i\,X^{-1}=Y\,A^i\,Y^{-1}$
    for all~$i$, then there exists a nonzero scalar
    $u \in \C^{\times}$ such that $Y=u\cdot X$.
\end{lemma}

\begin{proof}\leanok\uses{def:injective, lem:center_scalar}
    The matrix $Z=Y^{-1}X$ commutes with every $A^i$.
    Since $A$ is injective, $\{A^i\}$ spans $\MN{D}$, so $Z$
    lies in the centre of $\MN{D}$.
    By Lemma~\ref{lem:center_scalar}, $Z$ is a scalar matrix.
    Since $Z$ is invertible, that scalar is nonzero, and
    $Y=u\cdot X$.
\end{proof}

\begin{theorem}[Gauge uniqueness from equality of MPV families]
    \label{thm:gauge_unique_up_to_scalar_sameMPV}
    \lean{MPSTensor.gauge_unique_up_to_scalar_of_sameMPV}
    \leanok
    \uses{def:injective, def:same_mpv}
    Let $A$ be an injective tensor and assume that $A$ and $B$ generate the
    same MPV family.
    If $X,Y \in \GL_{D}(\C)$ both satisfy
    $B^i=X\,A^i\,X^{-1}=Y\,A^i\,Y^{-1}$ for all $i$, then there is a
    nonzero scalar $u \in \C^{\times}$ such that $Y=u\cdot X$.
\end{theorem}

\begin{proof}
    \leanok
    \uses{lem:gauge_unique_up_to_scalar}
    This is exactly Lemma~\ref{lem:gauge_unique_up_to_scalar}; the extra MPV
    hypothesis specifies the setting in which the two gauges are usually
    produced.
\end{proof}

\begin{definition}[Scalar 2-cochain]\label{def:scalar_cocycle}
    \lean{TNLean.Algebra.ScalarCocycle}
    \leanok
    A \emph{scalar 2-cochain} on a group $G$ is a function
    $\omega : G \times G \to \C^{\times}$.
\end{definition}

\begin{definition}[Projective representation]\label{def:projective_rep}
    \lean{TNLean.Algebra.ProjectiveRepresentation}
    \leanok
    \uses{def:scalar_cocycle}
    Given a scalar 2-cochain $\omega$, a \emph{projective representation}
    of $G$ on $\C^D$ is a map $\rho : G \to \GL_{D}(\C)$ satisfying, for all
    $g,h \in G$,
    \begin{align}
        \rho(g) \rho(h)
        &=\omega(g,h) \rho(gh).
        \label{eq:symmetry_projective_law}
    \end{align}
\end{definition}

\begin{lemma}[Associativity forces the cocycle condition]%
    \label{lem:cocycle_of_assoc}
    \lean{TNLean.Algebra.ProjectiveRepresentation.cocycle_of_assoc}
    \leanok
    \uses{def:projective_rep, def:scalar_cocycle}
    If $\rho$ is a projective representation with factor system $\omega$
    and $D > 0$, then $\omega$ satisfies the 2-cocycle identity
    \begin{align}
        \omega(g,h) \omega(gh,k)
        &=\omega(g,hk) \omega(h,k).
        \label{eq:symmetry_cocycle_identity}
    \end{align}
\end{lemma}

\begin{proof}\leanok\uses{def:projective_rep}
    Expand $\rho(g)\,(\rho(h)\rho(k))$ and
    $(\rho(g)\rho(h))\rho(k)$ using the projective multiplication law
    (\ref{eq:symmetry_projective_law}) and equate the results via matrix
    associativity. The hypothesis $D > 0$ allows cancellation of the
    invertible matrix factor, giving
    (\ref{eq:symmetry_cocycle_identity}).
\end{proof}

\begin{theorem}[Virtual representation theorem for injective MPS]%
    \label{thm:virtual_rep_injective}
    \lean{MPSTensor.virtual_rep_of_symmetric_injective}
    \leanok
    \uses{def:injective, def:on_site_symmetric, def:projective_rep,
        def:scalar_cocycle, lem:gauge_unique_up_to_scalar,
        lem:twisted_tensor_mul}
    Let $A$ be an injective MPS tensor that is on-site symmetric
    under a group representation $U : G \to \GL_{d}(\C)$.
    Then there exist:
    \begin{enumerate}
        \item a scalar 2-cochain
              $\omega : G \times G \to \C^{\times}$, and
        \item a map $\rho : G \to \GL_{D}(\C)$ satisfying, for all $g,h \in G$,
              \begin{align}
                  \rho(g) \rho(h)
                  &=\omega(g,h) \rho(gh).
                  \notag
              \end{align}
    \end{enumerate}
    such that, defining the gauge matrices $X(g):=\rho(g^{-1})$,
    one has, for every $g \in G$ and $i \in \{0,\ldots,d{-}1\}$,
    \begin{align}
        \widetilde{A}_g^i
        &=X(g)\,A^i\,X(g)^{-1}.
        \label{eq:symmetry_virtual_gauge}
    \end{align}
    In tensor-network notation, for each fixed~$g$ this is the local gauge
    relation
    \begin{align}
        &
        % Formula/source: Papers/0802.0447/StringOrder-v10.tex,
        % lines 328--349, gives the unitary u,V specialization
        % (V^\dagger = V^{-1}); the surrounding fundamental theorem supplies
        % the analogous general GL gauge relation displayed here.
        % Ink-to-index map: i labels the free physical leg and the west/east
        % wires carry the local tensor's two virtual indices.
        % Contraction set: the two internal horizontal gauge bonds are summed.
        % Type verdict: a rank-three local tensor for each fixed g.
        % Boundary signature: virtual-west=1 | virtual-east=1 |
        % physical-up=1 | physical-down=0.
        \begin{tenkz}[physical=up]
            \tnX{X(g)} & \tn[up=$i$]{A} & \tnX{X(g)^{-1}}
        \end{tenkz}
        \notag
    \end{align}
    with the black node denoting the tensor named in the surrounding
    formula and the two red side nodes acting on the virtual legs.
    In particular, $\rho$ is a projective representation of $G$ on the
    bond space $\C^D$.
\end{theorem}

\begin{proof}
    \leanok
    \uses{def:on_site_symmetric, lem:gauge_equiv_twisted,
        lem:gauge_unique_up_to_scalar, lem:twisted_tensor_mul}
    Since $A$ is injective and on-site symmetric, each twisted tensor
    $\widetilde{A}_g$ is gauge equivalent to~$A$
    (Lemma~\ref{lem:gauge_equiv_twisted}).
    Choose $X(g) \in \GL_{D}(\C)$ satisfying
    (\ref{eq:symmetry_virtual_gauge}) for all~$i$.

    By functoriality (Lemma~\ref{lem:twisted_tensor_mul}), the product
    $X(h)\,X(g)$ also intertwines $A$ with $\widetilde{A}_{gh}$.
    Since $X(gh)$ does the same, gauge uniqueness
    (Lemma~\ref{lem:gauge_unique_up_to_scalar}) gives a nonzero scalar
    $\omega'(h,g)$ such that
    \begin{align}
        X(h)\,X(g)
        &=\omega'(h,g)\,X(gh).
        \label{eq:symmetry_gauge_factor}
    \end{align}

    Defining $\rho(g):=X(g^{-1})$ converts
    (\ref{eq:symmetry_gauge_factor}) to the standard law
    (\ref{eq:symmetry_projective_law}), where
    $\omega(g,h):=\omega'(g^{-1},h^{-1})$.
\end{proof}

\begin{corollary}[2-cocycle identity for the factor system]%
    \label{cor:virtual_rep_cocycle}
    \lean{MPSTensor.virtual_rep_of_symmetric_injective_cocycle}
    \leanok
    \uses{thm:virtual_rep_injective, lem:cocycle_of_assoc}
    If $D > 0$, the factor system $\omega$ from
    Theorem~\ref{thm:virtual_rep_injective} satisfies the
    2-cocycle identity for all $g,h,k \in G$:
    \begin{align}
        \omega(g,h) \omega(gh,k)
        &=\omega(g,hk) \omega(h,k).
        \notag
    \end{align}
\end{corollary}

\begin{proof}
    \leanok
    \uses{thm:virtual_rep_injective, lem:cocycle_of_assoc}
    Theorem~\ref{thm:virtual_rep_injective} gives $\rho$ and $\omega$.
    By Lemma~\ref{lem:cocycle_of_assoc}, the cocycle identity
    (\ref{eq:symmetry_cocycle_identity}) follows from associativity of
    matrix multiplication and invertibility of the representation matrices.
\end{proof}

\section{Cohomology Classes of Virtual Symmetry Cocycles}

Two cocycles may differ by a coboundary yet represent the same
projective-representation class. The following definitions and results
make this precise and establish the quotient $H^2(G,\C^{\times})$.

\begin{definition}[Coboundary]\label{def:is_coboundary}
    \lean{TNLean.Algebra.ScalarCocycle.IsCoboundary}
    \leanok
    \uses{def:scalar_cocycle}
    A scalar 2-cocycle $\omega$ is a \emph{coboundary} if there exists
    $\varphi \colon G \to \C^{\times}$ such that, for all $g,h \in G$,
    \begin{align}
        \omega(g,h)
        &=\varphi(g)\varphi(h)\varphi(gh)^{-1}.
        \label{eq:symmetry_coboundary}
    \end{align}
\end{definition}

\begin{definition}[Cohomologous cocycles]\label{def:cohomologous_to}
    \lean{TNLean.Algebra.ScalarCocycle.CohomologousTo}
    \leanok
    \uses{def:scalar_cocycle}
    Two cocycles $\omega_1,\omega_2$ are \emph{cohomologous}
    (written $\omega_1 \sim \omega_2$) if there exists
    $\varphi \colon G \to \C^{\times}$ such that, for all $g,h \in G$,
    \begin{align}
        \omega_1(g,h)
        &=\varphi(g)\varphi(h)\varphi(gh)^{-1}\omega_2(g,h).
        \label{eq:symmetry_cohomologous}
    \end{align}
\end{definition}

\begin{theorem}[Cohomology relation is an equivalence relation]%
    \label{thm:cohomologous_equivalence}
    \lean{TNLean.Algebra.ScalarCocycle.CohomologousTo.equivalence}
    \leanok
    \uses{def:cohomologous_to}
    The relation ``cohomologous to'' is reflexive, symmetric, and
    transitive on the set of scalar 2-cocycles.
\end{theorem}

\begin{proof}\leanok\uses{def:cohomologous_to}
    Reflexivity: take $\varphi=1$.
    Symmetry: replace $\varphi$ by $\varphi^{-1}$.
    Transitivity: compose the witnesses $\varphi$ and $\psi$ pointwise
    in (\ref{eq:symmetry_cohomologous}).
\end{proof}

\begin{lemma}[Coboundary iff cohomologous to the trivial cocycle]%
    \label{lem:coboundary_iff_cohomologous_one}
    \lean{TNLean.Algebra.ScalarCocycle.isCoboundary_iff_cohomologousTo_one}
    \leanok
    \uses{def:is_coboundary, def:cohomologous_to}
    A cocycle $\omega$ is a coboundary if and only if
    $\omega \sim \Id$, where $\Id(g,h)=1$ for all $g,h$.
\end{lemma}

\begin{proof}\leanok\uses{def:is_coboundary, def:cohomologous_to}
    Both directions follow by absorbing or introducing the factor
    $\Id(g,h)=1$ via $x\cdot 1=x$ in
    (\ref{eq:symmetry_cohomologous}).
\end{proof}

\begin{theorem}[Gauge independence of the cocycle class]%
    \label{thm:cocycle_gauge_independence}
    \lean{MPSTensor.cohomologousTo_of_isInjective}
    \leanok
    \uses{thm:virtual_rep_injective, lem:gauge_unique_up_to_scalar,
        def:cohomologous_to}
    If the bond dimension satisfies $D \ge 1$ and two projective
    representations $\rho_1,\rho_2$ (with cocycles
    $\omega_1,\omega_2$) both arise as virtual representations of the
    same injective MPS tensor $A$ under an on-site symmetry~$U$, then
    $\omega_1 \sim \omega_2$.
\end{theorem}

\begin{proof}\leanok
    \uses{thm:virtual_rep_injective, lem:gauge_unique_up_to_scalar,
        def:cohomologous_to}
    For each $g$, gauge uniqueness gives a nonzero scalar $f(g)$ with
    $X_2(g^{-1})=f(g)\,X_1(g^{-1})$.
    Equivalently, for each fixed $g$ the two local gauges
    \begin{tenkzequation}
        % Formula/source: Papers/0802.0447/StringOrder-v10.tex,
        % lines 328--349, gives the unitary u,V specialization
        % (V^\dagger = V^{-1}); the surrounding fundamental theorem supplies
        % the analogous general GL gauge relation used here.
        % Ink-to-index map: i is free and the outer horizontal wires carry
        % the west/east virtual indices for the first gauge choice.
        % Contraction set: the two internal horizontal bonds are summed.
        % Type verdict: a rank-three local tensor.
        % Boundary signature: virtual-west=1 | virtual-east=1 |
        % physical-up=1 | physical-down=0.
        \begin{tenkz}[physical=up]
            \tnX{X_1(g^{-1})} & \tn[up=$i$]{A} &
            \tnX{X_1(g^{-1})^{-1}}
        \end{tenkz}
        \qquad \text{and} \qquad
        % Formula/source: the same source passage, with the second virtual
        % gauge X_2(g^{-1}) in place of X_1(g^{-1}).
        % Ink-to-index map: i is free; the outer wires are the two virtual
        % indices for the second gauge choice.
        % Contraction set: the two internal horizontal bonds are summed.
        % Type verdict: a rank-three local tensor, comparable to the left picture.
        % Boundary signature: virtual-west=1 | virtual-east=1 |
        % physical-up=1 | physical-down=0.
        \begin{tenkz}[physical=up]
            \tnX{X_2(g^{-1})} & \tn[up=$i$]{A} &
            \tnX{X_2(g^{-1})^{-1}}
        \end{tenkz}
    \end{tenkzequation}
    describe the same twisted tensor, so they differ by a scalar.
    Substituting into the projective law
    (\ref{eq:symmetry_projective_law}) and cancelling $X_1(g\cdot h)$
    yields
    \begin{align}
        \omega_2(g,h)
        &=f(g)\,f(h)\,f(gh)^{-1}\omega_1(g,h).
        \label{eq:symmetry_gauge_coboundary}
    \end{align}
    Thus, (\ref{eq:symmetry_gauge_coboundary}) has precisely the form
    required by (\ref{eq:symmetry_cohomologous}).
\end{proof}

\begin{definition}[Multiplicative 2-cocycle condition]\label{def:is_cocycle}
    \lean{TNLean.Algebra.ScalarCocycle.IsCocycle}
    \leanok
    A scalar 2-cochain $\omega \colon G \times G \to \C^{\times}$ is a
    \emph{2-cocycle} if it satisfies, for all $g,h,k \in G$,
    \begin{align}
        \omega(g,h) \omega(gh,k)
        &=\omega(g,hk) \omega(h,k).
        \notag
    \end{align}
\end{definition}

\begin{definition}[Second cohomology quotient]\label{def:h2_cx}
    \lean{TNLean.Algebra.H2}
    \leanok
    \uses{def:cohomologous_to, def:is_cocycle}
    The second cohomology set is the quotient
    \begin{align}
        H^2(G,\C^{\times})
        &:=Z^2(G,\C^{\times})/{\sim},
        \label{eq:symmetry_h2_quotient}
    \end{align}
    where $\sim$ is the coboundary-equivalence relation.
\end{definition}

\begin{definition}[Projective equivalence]\label{def:projectively_equivalent}
    \lean{TNLean.Algebra.ProjectivelyEquivalent}
    \leanok
    \uses{def:projective_rep, def:cohomologous_to}
    Two projective representations are \emph{projectively equivalent}
    when their factor systems are cohomologous.
\end{definition}

\begin{theorem}[Projective equivalence and cohomologous cocycles]%
    \label{thm:projrep_equiv_iff_cohomologous}
    \lean{TNLean.Algebra.projRep_equiv_iff_cohomologous}
    \leanok
    \uses{def:projectively_equivalent, def:cohomologous_to}
    Let $\rho_1,\rho_2$ be projective representations with factor systems
    $\omega_1,\omega_2$. Then
    \begin{align}
        \rho_1 \sim_{\mathrm{proj}} \rho_2
        &\iff \omega_1 \sim \omega_2.
        \label{eq:symmetry_projective_equivalence}
    \end{align}
\end{theorem}

\begin{proof}\leanok\uses{def:projectively_equivalent}
    This follows from the definition of projective equivalence as
    cohomology-class equality for factor systems.
\end{proof}

\subsection{Non-triviality via the commutator phase}

For an abelian symmetry group the cohomology class of a factor system is
detected by an antisymmetric pairing, which gives a computable test for
non-triviality.

\begin{definition}[Commutator phase]\label{def:comm_phase}
    \lean{TNLean.Algebra.ScalarCocycle.commPhase}
    \leanok
    \uses{def:scalar_cocycle}
    The \emph{commutator phase} of a $2$-cochain $\omega$ at a pair $g,h$ is
    $\beta_\omega(g,h):=\omega(g,h) \omega(h,g)^{-1}$. It is defined for any
    $\C^{\times}$-valued $2$-cochain; the cocycle condition is not needed.
\end{definition}

\begin{theorem}[The commutator phase is a class invariant]\label{thm:comm_phase_invariant}
    \lean{TNLean.Algebra.ScalarCocycle.commPhase_eq_of_cohomologousTo}
    \leanok
    \uses{def:comm_phase, def:cohomologous_to}
    If $\omega_1\sim\omega_2$ and $gh=hg$, then
    $\beta_{\omega_1}(g,h)=\beta_{\omega_2}(g,h)$.
\end{theorem}

\begin{proof}\leanok\uses{def:comm_phase, def:cohomologous_to}
    Write
    $\omega_1(g,h)=\varphi(g)\varphi(h)\varphi(gh)^{-1}\omega_2(g,h)$
    for a coboundary witness $\varphi$. Substituting into the ratio and using
    $gh=hg$ and the commutativity of $\C^{\times}$ gives
    \begin{align}
        \beta_{\omega_1}(g,h)
        &=
        \frac{\varphi(g)\varphi(h)\varphi(gh)^{-1} \omega_2(g,h)}
             {\varphi(h)\varphi(g)\varphi(hg)^{-1} \omega_2(h,g)}
        \notag\\
        &=\frac{\omega_2(g,h)}{\omega_2(h,g)}
        =\beta_{\omega_2}(g,h),
        \label{eq:symmetry_comm_phase_invariance}
    \end{align}
    where the $\varphi$-factors cancel since
    $\varphi(gh)=\varphi(hg)$.
\end{proof}

\begin{definition}[Non-trivial cohomology class]\label{def:nontrivial_class}
    \lean{TNLean.Algebra.ScalarCocycle.IsNontrivialClass}
    \leanok
    \uses{def:cohomologous_to}
    A factor system $\omega$ has a \emph{non-trivial class} when it is not
    cohomologous to the constant cocycle $1$.
\end{definition}

\begin{theorem}[Commutator-phase non-triviality test]\label{thm:nontrivial_of_comm_phase}
    \lean{TNLean.Algebra.ScalarCocycle.isNontrivialClass_of_commPhase_ne_one}
    \leanok
    \uses{def:nontrivial_class, def:comm_phase, thm:comm_phase_invariant}
    If a commuting pair $g,h$ has $\beta_\omega(g,h)\neq 1$, then $\omega$ has a
    non-trivial class.
\end{theorem}

\begin{proof}\leanok\uses{thm:comm_phase_invariant}
    The trivial cocycle has commutator phase $1$ at every pair, so by
    Theorem~\ref{thm:comm_phase_invariant} and
    (\ref{eq:symmetry_comm_phase_invariance}), any cocycle cohomologous to it
    has $\beta=1$ at every commuting pair. Contrapositively,
    $\beta_\omega(g,h)\neq 1$ for one commuting pair forces $\omega$ to be
    non-cohomologous to the trivial cocycle.
\end{proof}
